% Part: incompleteness
% Chapter: incompleteness-provability
% Section: godels-paper

\documentclass[../../../include/open-logic-section]{subfiles}

\begin{document}

\olfileid[tr]{inc}{inp}{gop}

\olsection{G\"odel'in Özgün Makalesiyle Karşılaştırma}

G\"odel'in 1931 tarihli makalesine biraz zaman ayırmak yararlı olur. Giriş
bölümü, az önce tartıştığımız fikirlerin bir taslağını verir. Makaleyi yalnızca
gözden geçirseniz bile her aşamada ne yapıldığını görmek kolaydır: G\"odel
önce $P$ biçimsel sistemini (sözdizimini, aksiyomları ve ispat kurallarını)
betimler; sonra ilkel özyinelemeli fonksiyonları ve bağıntıları tanımlar;
ardından $x B y$'nin ilkel özyinelemeli olduğunu gösterir ve ilkel
özyinelemeli fonksiyonlarla bağıntıların $\Th{P}$ içinde temsil edildiğini
savunur. Sonra yukarıdaki gibi eksiklik teoremini ispatlar. 3. Bölümde,
ispatlanamayan önermenin aritmetik dilinde bir tümce olarak alınabileceğini
gösterir. Dizileri ele alıp özyinelemeli fonksiyonların $\Th{Q}$ içinde temsil
edilebilir olduğunu göstermek için bizim de kullandığımız $\beta$-lemmasının
kökeni budur. G\"odel, yeterli olan en küçük aksiyomlar kümesini ayıracak kadar
ileri gitmez; fakat artık $\Th{Q}$'nun bu işi gördüğünü biliyoruz. Son olarak
4. Bölümde ikinci eksiklik teoreminin ispatını ana hatlarıyla verir.

\end{document}
